\begin{definition}[完全正可分动力学]
对固定、与系统无关的环境初态所产生的一族 CPTP 映射 $\Phi_{t:0}$，
若任意 $t\ge s\ge0$ 都存在 CPTP 映射 $V_{t:s}$，使
$\Phi_{t:0}=V_{t:s}\Phi_{s:0}$，则称这族动力学完全正可分。
定义不要求 $\Phi_{s:0}$ 可逆；可逆时中间映射唯一为
$V_{t:s}=\Phi_{t:0}\Phi_{s:0}^{-1}$。
\end{definition}

\begin{proposition}[可逆时间局域方程的瞬时率判据]
在有限维工作空间内，设 $\Phi_{t:0}$ 可微且可逆。
若其生成元在正交无迹算符基中取标准形式
\[
\mathcal K_t(\rho)=-\frac{\ii}{\hbar}[H_t,\rho]
+\sum_a\gamma_a(t)
\left(L_a\rho L_a^\dagger-\frac12\{L_a^\dagger L_a,\rho\}\right),
\]
其中 $\gamma_a$ 是耗散系数矩阵的本征值，那么在没有奇点的区间上，
动力学完全正可分当且仅当所有 $\gamma_a(t)\ge0$。
\end{proposition}
\begin{proof}
非负率使每个短时间生成元具有 GKSL 形式。
把短时 CPTP 传播子按时间相乘并取极限，得到任意 $V_{t:s}$ 的完全正保迹性。
反向考察 $V_{t+\dd t:t}=I+\dd t\,\mathcal K_t+o(\dd t)$ 的 Choi 算符。
在恒等信道的 Choi 向量正交补上，其零阶部分为零；
Hamiltonian 和反对易子项在两侧同时投影后也为零。
剩下的一阶二次型恰是耗散系数矩阵乘 $\dd t$。
完全正要求这一正交补上的二次型非负，因此该矩阵的标准本征率必须非负。
一般非正交的跳跃算符写法并不唯一，不能直接拿任意一组展开系数充当这个判据。
\end{proof}

\begin{exbox}[有限环境中的相干回流与 TCL 奇点]
以两条电子通道作系统 $S$，保留一个二态环境 $B$。
初态因子化，环境取 $I_B/2$，相互作用绘景 Hamiltonian 为
$H_{SB}=\hbar\Omega\sigma_z^{(S)}\otimes\sigma_z^{(B)}/2$。
环境两个本征态分别让系统转过相反的相位，所以
\begin{align*}
\Phi_t(\rho)&=\frac12U_+(t)\rho U_+(t)^\dagger
+\frac12U_-(t)\rho U_-(t)^\dagger,\\
U_\pm(t)&=\exp(\mp\ii\Omega t\sigma_z/2),\\
\rho_{01}(t)&=\nu(t)\rho_{01}(0),\qquad\nu(t)=\cos(\Omega t).
\end{align*}
这是两个幺正信道的凸组合，故每一时刻都完全正。
在 $\nu(t)\ne0$ 的区间，时间局域方程是
\begin{equation}
\dot\rho=\gamma(t)\mathcal D[\sigma_z]\rho,
\qquad\gamma(t)=-\frac{\dot\nu(t)}{2\nu(t)}
=\frac{\Omega}{2}\tan(\Omega t).
\label{eq:feqo-exact-finite-memory-dephasing-revival}
\end{equation}
当 $\pi/2<\Omega t<\pi$ 时 $\gamma(t)<0$，相干模长回升。
这个负率并没有使任何从初始时刻传播的态变成负态；
它表示从某个中间时刻任意重新制备的系统态，不能忽略当时已经形成的环境相关来传播。

取两个初态 $\rho_\pm(0)=(I\pm\sigma_x)/2$，其迹距离为
$\tfrac12\|\Phi_t(\rho_+)-\Phi_t(\rho_-)\|_1=|\cos(\Omega t)|$。
距离的回升直接显示之前流入环境的可区分性返回系统。
在 $t_*=\pi/(2\Omega)$，两个系统态完全相同，但随后又变得不同；
任何只作用于这个相同中间态的映射都不可能产生两种不同的输出。
故跨越 $t_*$ 的中间 CPTP 映射不存在。
同时 $\Phi_t$ 在 $t_*$ 不可逆，TCL 率发散，而实际密度矩阵仍有界且光滑。
\end{exbox}
